\section{Introduction}
\label{sec:slam-intro}
This chapter focuses on the development and implementation of a SLAM system designed to operate independently of frequent GNSS updates. The primary objective is to ensure high-fidelity localisation in GNSS-denied or GNSS-degraded scenarios. Furthermore, a critical requirement for this system is temporal robustness; that is, the SLAM front-end must remain invariant to fluctuating visual appearances, enabling continuous 24/7 operation across both diurnal and nocturnal cycles, different weather conditions, and plant growth cycles. Thirdly, as aforementioned, vineyard and orchard rows form long, symmetric, corridor-like structures which limits the geometric diversity available to scan-matching algorithms, making them susceptible to ambiguity and longitudinal drift. The system presented in this chapter is therefore not intended as a general-purpose SLAM solution; it is a specialised architecture targeted at the structural and perceptual characteristics of row-crop environments such as vineyards and orchards.

\subsection{Motivation}
\label{subsec:slam-motivation}
The necessity for this research stems from the operational limitations of the current UCVision Rover localisation system. Prior iterations of this platform utilised an Extended Kalman Filter (EKF) heavily reliant on Real-Time Kinematic (RTK) GPS data. While precise under open skies, this dependency proved non-viable for robust field operations where tree canopies, topography, and public RTK base station blackouts frequently cause signal multipath or dropout. Adverse weather further compounds this unreliability: heavy precipitation increases tropospheric delay errors \cite{solheimPropagationDelaysGPSSignals1999} and multipath reflections from wet vegetation \cite{nievinskiForwardModelingGPSMultipath2014}, while ionospheric disturbances — which are more frequent during solar weather events and magnetic storms — introduce rapid carrier-phase cycle slips that break the integer ambiguity fix and degrade the RTK solution from centimetre-level fixed accuracy to decimetre-level float accuracy or worse \cite{wielgoszAnalysisNetworkRTKIonospheric2005}. In practice, a significant proportion of field scans conducted with the UCVision Rover were rendered unusable due to the inability to acquire or sustain a consistent RTK fix, causing the localisation system to fail entirely.

Transitioning to a LiDAR-centric SLAM approach, tailored to these agricultural environments, provides three distinct advantages:
\begin{itemize}
	\item \textbf{GNSS-independent localisation:} By shifting the primary motion estimation to scan-matching and loop closure, the system becomes less reliant on RTK.
	\item \textbf{Corridor robustness:} The architecture presented here specifically addresses the structural challenges of row-crop environments through vineyard-focused parameterisation and sensor fusion suited to row-structured terrain.
	\item \textbf{Downstream utility:} Beyond localisation, the high-density LiDAR maps generated can be used for additional tasks such as landmark detection (e.g., vine trunks and trellis posts), spatial alignment of Gaussian splats captured from opposing sides of a row, and global registration of data across multiple vineyard bays.
\end{itemize}
